\section{综合：名分、宫门与区域中心}

汉的皇帝名义有实际效力：它能任命、赐爵、组织仪式，也能让竞争者在诏令中表达自己。它却从不自动制造粮食、军队或地方服从。吕后临朝、霍光辅政、王莽居摄、邓太后与窦氏、何进和宦官的冲突，反复显示继承不是血缘的自然延续，而是禁军、诏令、外戚、官僚和可接受名分的重新组合。

两次受禅尤其适合比较。九年新朝受禅和二二〇年魏受禅都保留了符命、劝进或诏册一类仪式语言；它们并非可忽略的虚饰，因为仪式把权力重组置入可被官署书写和承认的形式。但仪式不能单独证明自愿、共识或全国控制。前者之前有王氏对辅政与人事的占据，后者之前有曹氏对北方军政、屯田和朝廷的控制；二者都说明名分必须搭在先已形成的资源关系上。

汉末诸中心不应被写作几位人物在空白地图上的分割。冀州、许县、荆州、益州、江东和关中各有仓储、河道、山口、宗族和军事网络。官渡把河北与河南资源更紧地接入曹氏，赤壁使长江的条件显出；这些胜负改变了区域间的选择，并没有立即消除局部社会的战争、迁徙和征发。魏、蜀、吴后来各自继承汉代的官名、文书或名分，恰好证明帝国的遗产是可争夺的接口，而不是一件可完整转让的器物。
